% ==============================================================================
% sections/results.tex
% Created: 2026-05-08
% ==============================================================================

\subsection{Baseline estimates}\label{subsec:results:baseline}

Table~\ref{tab:baseline_implied} reports implied entry and exit probabilities
from the baseline model, and Table~\ref{tab:baseline_params} reports the
underlying structural parameter estimates. The columns cover the six
specification variants: stationarity-restricted and free initial probability,
crossed with no misclassification, symmetric misclassification, and asymmetric
misclassification.

Without misclassification correction---the standard Markov transition matrix
applied to observed employment status---the free-initial-condition estimates
of the quarterly entry and exit rates are 8.28\% (SE 0.05) and 8.54\%
(SE 0.05). Imposing stationarity gives 8.18\% and 8.64\%, respectively. These estimates are close to the raw
transition rates in Table~\ref{tab:descriptive}, confirming that the
no-misclassification specification simply recovers the moments of the observed
transition matrix.

The symmetric misclassification specification yields a dramatically different
picture. The estimated misclassification probability is $\hat{\pi} = 2.99\%$
in both initial-condition specifications. The stationary model implies entry
and exit rates of 2.86\% (SE 0.05) and 3.04\% (SE 0.05); the free-initial-condition
model gives 2.98\% (SE 0.06) and 2.91\% (SE 0.06).
The correction factor is approximately three: a single additional parameter
reduces the estimated transition rates by two-thirds. This is the main finding
of the paper. The factor-of-three correction is driven by the severity of the
Markov violation documented in Section~\ref{sec:data:puzzle}: even a small
misclassification probability inflates high-frequency transition rates
substantially because each transient error fabricates two transitions.

The asymmetric misclassification specification, which allows the probabilities
of falsely reporting employment and non-employment to differ, produces estimates
that are virtually indistinguishable from the symmetric specification. The
estimated false-positive and false-negative probabilities are approximately
3.05\% and 2.93\%, and the implied transition rates are within a rounding error
of the symmetric case. The directional difference is only about 0.12 percentage
points, although it is statistically detectable in this large sample. We adopt
the symmetric specification subsequently because the economic implications are
unchanged and the restriction is substantially more parsimonious.

Comparing the stationarity-restricted and free initial-probability specifications
yields modest differences in the transition rates but a larger difference in
the implied steady-state employment share. We therefore report both versions
for the baseline model and use a free initial probability in the revised
covariate comparison, where transition-varying characteristics also make a
time-homogeneous stationarity restriction unattractive.

The analytical survey-weighted sandwich standard errors indicate that the
misclassification probability is estimated with considerable precision. The
delta-method 95\% interval for $\hat{\pi}$ in the symmetric
stationarity-restricted specification is approximately [2.92\%, 3.06\%],
comfortably above zero and well below magnitudes that would
be necessary to explain the raw transition matrix moments without correction.

\subsection{Observable worker heterogeneity}\label{subsec:results:covariates}

A natural concern is that the estimated misclassification probability is
capturing heterogeneity in workers' true transition rates rather than
measurement error. Workers who are persistently close to the employment margin
may exhibit frequent transitions in the data even if the Markov assumption holds;
if these workers are a small share of the sample, their high transition rates
could generate Markov violations when the average transition rate is treated as
homogeneous.

Table~\ref{tab:cov_risk_weighted} reports the implied transition rates from the
observable heterogeneity extension. We consider three covariate sets of
increasing richness: (1) age and education only; (2) additionally controlling
for gender and race; and (3) additionally controlling for origin-wave contract
type and informal-sector employment in the exit equation. Rural or urban
residence is not used, so the sample is not restricted on its availability.
All columns estimate the initial employment probability freely. Rates are
averaged with survey weights and posterior probabilities of membership in the
relevant origin-state risk set.

In the no-error models, entry is 8.14\% and exit is 8.58\% per quarter in all
three sets. With symmetric reporting error, the corresponding rates range from
3.62--3.99\% and 3.71--4.16\%, while the estimated per-wave error probability
declines only modestly from 2.58\% to 2.39\%. Expected entry and exit flows are
nearly balanced in every specification; Table~\ref{tab:cov_risk_weighted}
therefore reports their sum as total churn, rather than presenting the two
component flows separately. The appendix shows substantial
heterogeneity: in symmetric Set 3, the risk-weighted 10th--90th percentile
ranges are 1.07--7.97\% for entry and 0.46--9.61\% for exit. A permanent
contract raises employment persistence, whereas informal-sector employment
lowers it. Thus observable heterogeneity changes who faces high transition
risk, but does not eliminate the estimated reporting-error component; ignoring
it continues to overstate average transition hazards by roughly a factor of two.

\subsection{Unobserved worker heterogeneity}\label{subsec:results:fmm}

The finite mixture model (FMM) extension relaxes the assumption that workers
form a homogeneous population with respect to their underlying transition
probabilities. Table~\ref{tab:fmm} reports the FMM results.

Without accounting for misclassification, the two-type FMM assigns approximately
equal weight to each type ($\hat{\phi} \approx 0.50$). Type B has very high
estimated entry and exit rates ($\hat{\theta}_0^B \approx 0.37$,
$\hat{\theta}_1^B \approx 0.57$), while type A has much lower rates. This
``coin-flip'' type B is an implausible description of any group of actual workers;
it is more naturally interpreted as a misclassification artefact, where a subset
of observations with particularly noisy employment signals is assigned to a
spurious type.

Adding misclassification to the FMM substantially alters this picture. The mixing
proportion shifts to $\hat{\phi} = 43.26\%$, the type A entry and exit rates
fall to 6.05\% and 1.95\%, and the type B rates become 2.13\% and 5.44\%.
These are plausible transition rates for two distinct worker groups---perhaps
corresponding to workers with stable versus unstable employment relationships.
The estimated misclassification probability in the FMM with $\pi$ is
$\hat{\pi} = 2.92\%$, virtually identical to the baseline. This is the key
result: allowing the model to explain Markov violations through either
unobserved heterogeneity or misclassification, it attributes approximately the
same fraction to misclassification as the simpler homogeneous model. Unobserved
heterogeneity does not absorb the misclassification signal; instead, adding
misclassification resolves the implausible coin-flip type.

\subsection{Inconsistency-augmented misclassification}\label{subsec:results:inconsistency}

Table~\ref{tab:inconsistency} reports the results of the inconsistency-augmented
model. Both response error and incorrect panel linkage can generate age or
education inconsistencies alongside employment-status misclassification.
Conditioning misclassification on these indicators therefore tests whether
unreliable linked records contain more misclassification-consistent employment
histories, without attempting to distinguish the two mechanisms.

In the preferred free-initial-probability specification, the baseline
misclassification probability for a wave with no attributed age or education
inconsistency is 2.26\%. An age inconsistency raises the estimated probability
by 15.74 percentage points and an education inconsistency raises it by 4.24
percentage points. The survey-weighted mean misclassification probability is
3.35\%. All three estimates are reported with analytical standard errors in
Table~\ref{tab:inconsistency}.

These estimates show that unreliable linked records are substantially more
likely to exhibit employment histories that the model attributes to
misclassification. They do not distinguish response error from incorrect panel
matches: either mechanism can jointly generate inconsistent demographic
responses and implausible employment sequences. The association remains large
when the true entry and exit processes are allowed to differ between records
with and without any inconsistency (Table~\ref{tab:inconsistency_reliability}).
It is also increasing in severity for age inconsistencies, while the severity
gradient is much smaller for education (Table~\ref{tab:inconsistency_extent}).
